% =========================================================
% Control: Continuity — Notes Index
% =========================================================

% ---------------------------------------------------------
% Breadcrumb
% ---------------------------------------------------------
\begin{tcolorbox}[
  colback=gray!6,
  colframe=gray!40,
  arc=2pt,
  left=8pt, right=8pt, top=6pt, bottom=6pt,
  title={\small\textbf{Where You Are in the Journey}},
  fonttitle=\small\bfseries
]
\begin{center}
\small
$\mathbb{N}$, $\mathbb{Z}$, $\mathbb{Q}$, $\mathbb{R}$
$\;\to\;$ Real Analysis
$\;\to\;$ Distance (Metric Spaces)
$\;\to\;$ Nearness (Topology)
$\;\to\;$ Finiteness (Compactness)
$\;\to\;$ \textbf{Control (Continuity)}
$\;\to\;$ Accumulation (Integration)
$\;\to\;$ $\cdots$
\end{center}

\medskip
\noindent\textbf{How we got here.}
Topology gave us the language of open sets.
Compactness gave us finite control over infinite covers.
Continuity is what connects these structural tools to the
behavior of functions: it is the condition that a map
\emph{respects} the neighborhood structure built in the
previous two chapters.

\medskip
\noindent\textbf{What this chapter builds.}
The $\varepsilon$-$\delta$ definition of continuity at a point,
the four equivalent characterizations (metric, topological,
sequential, and limit-based), continuity on a set,
the Extreme Value Theorem, the Intermediate Value Theorem,
uniform continuity, and the upgrade from pointwise to uniform
control that compact domains provide.

\medskip
\noindent\textbf{Where this leads.}
Uniform continuity on compact domains is the precise hypothesis
that makes Riemann sums partition-independent and forces
the Riemann integral to exist.
Differentiability is a strictly stronger condition: a function
is differentiable at a point only if it is continuous there,
but continuity does not imply differentiability.
\end{tcolorbox}
\vspace{1em}

% ---------------------------------------------------------
% Structural Role
% ---------------------------------------------------------
\begin{tcolorbox}[
  colback=gray!6,
  colframe=gray!40,
  arc=2pt,
  left=8pt, right=8pt, top=6pt, bottom=6pt,
  title={\small\textbf{Structural Role: Control — Preservation of Structure}},
  fonttitle=\small\bfseries
]
Control describes maps that respect neighborhood structure.

In metric language:
small input changes produce small output changes.

In topological language:
inverse images of open sets are open.

These two formulations are equivalent, and the equivalence is
not incidental --- it reveals that continuity is not about
formulas.
It is about structure preservation: a continuous map transports
convergence, compactness, and connectedness faithfully from
domain to codomain.

Once compactness is in place, continuity gains global power.
On compact sets, every continuous function is uniformly
continuous --- a single tolerance controls behavior across the
entire domain rather than varying point by point.

This upgrade from local to global control is the interaction
principle at the heart of analysis:
\[
  \text{Compactness} + \text{Continuity}
  \;\Longrightarrow\;
  \text{Global Stability.}
\]
Uniform continuity is one of the first deep rewards of the entire
structural progression, and the key that unlocks Riemann
integration in the next chapter.
\end{tcolorbox}
\vspace{1em}

% ---------------------------------------------------------
% Structural Roadmap
% ---------------------------------------------------------
\subsection*{Structural Roadmap}

\begin{center}
\textbf{Definitions $\longrightarrow$ Main Theorems
$\longrightarrow$ Consequences and Logical Structure}
\end{center}

The development follows the definition--theorem--structure
architecture, now applying the topological and compactness
machinery to functions.
The global progression is:

\begin{enumerate}
  \item Continuity at a point: $\varepsilon$-$\delta$ definition
  \item Four equivalent characterizations: metric, topological,
        sequential, limit-based
  \item Continuity on a set and on compact sets
  \item Extreme Value Theorem: compact domain $\Rightarrow$
        attains max and min
  \item Intermediate Value Theorem: connected domain $\Rightarrow$
        attains all intermediate values
  \item Uniform continuity: definition and comparison with
        pointwise continuity
  \item Uniform continuity on compact sets (Heine's theorem)
  \item Consequences for integration
\end{enumerate}

\begin{remark*}[Primary sources]
Abbott, \textit{Understanding Analysis}, Chapters~4--5;
Pugh, \textit{Real Mathematical Analysis}, Chapter~3;
Rudin, \textit{Principles of Mathematical Analysis}, Chapter~4.
\end{remark*}

% ---------------------------------------------------------
% Status
% ---------------------------------------------------------
\vspace{1em}
\begin{tcolorbox}[
  colback=gray!6, colframe=gray!40, arc=2pt,
  left=8pt, right=8pt, top=6pt, bottom=6pt,
  title={\small\textbf{Status: Planned}},
  fonttitle=\small\bfseries
]
\begin{center}\Large\bfseries Coming Soon\end{center}
\vspace{6pt}
\noindent Notes, proofs, and exercises will appear here in a future revision.
\end{tcolorbox}

% Future sections (uncomment as content is written):
% \input{volume-iii/analysis/continuity/notes/notes-definition}
% \input{volume-iii/analysis/continuity/notes/notes-characterizations}
% \input{volume-iii/analysis/continuity/notes/notes-evt-ivt}
% \input{volume-iii/analysis/continuity/notes/notes-uniform}
% \input{volume-iii/analysis/continuity/notes/notes-consequences}
